% LOCAL 2024/2025 — APMCM 亚太杯中文赛论文 LaTeX 模板 (基于 ctex article)
% 用法: 由 scripts/render_paper.py --competition apmcm 调用
% 编译: xelatex main.tex (三编, 中文必须 xelatex)
% APMCM 中文赛无统一公开模板, 本模板按中文数学建模优秀论文常见排版组织

\documentclass[12pt,a4paper]{ctexart}

\usepackage[margin=2.54cm]{geometry}
\usepackage{amsmath,amssymb,amsthm}
\usepackage{graphicx}
\usepackage{float}
\usepackage{booktabs}
\usepackage{longtable}
\usepackage{multirow}
\usepackage{listings}
\usepackage{xcolor}
\usepackage{hyperref}
\usepackage{url}
\usepackage{cite}
\usepackage{caption}
\usepackage{subcaption}
\usepackage{enumitem}
\usepackage{titlesec}
\usepackage{fancyhdr}
\usepackage{lastpage}

\hypersetup{
    colorlinks=true,
    linkcolor=blue,
    filecolor=magenta,
    urlcolor=blue,
    citecolor=blue,
}

\lstset{
    basicstyle=\small\ttfamily,
    breaklines=true,
    frame=single,
    numbers=left,
    numberstyle=\tiny,
    showstringspaces=false,
    extendedchars=false,
}

\pagestyle{fancy}
\fancyhf{}
\fancyhead[L]{APMCM 中文赛论文 \#XXXX}
\fancyhead[R]{第 \thepage\ 页 / 共 \pageref{LastPage} 页}
\renewcommand{\headrulewidth}{0.4pt}

\titleformat*{\section}{\Large\bfseries}
\titleformat*{\subsection}{\large\bfseries}
\titleformat*{\subsubsection}{\normalsize\bfseries}

\begin{document}

% ===== 封面 =====
\begin{titlepage}
\begin{center}
\vspace*{2cm}
{\Huge \textbf{APMCM 亚太杯中文赛论文标题}}\\[2cm]
{\Large 题号: \textbf{A}}\\[1cm]
{\Large 队伍编号: \textbf{XXXX}}\\[3cm]
\end{center}
\end{titlepage}

% ===== 摘要 =====
\section*{摘要}

% render_paper.py 注入锚: 下一行在渲染时被替换为 abstract 节内容; 手写论文时删除锚点行直接写
% __ABSTRACT__

\textbf{关键词:} 关键词1; 关键词2; 关键词3; 关键词4

\newpage
\tableofcontents
\newpage

% ===== 正文 =====
% render_paper.py 注入锚: 下方成对 SECTIONS 锚行之间的
% 演示正文会被 sections/*.tex 自动替换; 手写论文时删除两个锚点行直接写
% __SECTIONS__

\section{问题背景与重述}

\subsection{工程背景}

\subsection{题目重述}

\subsection{数据概览}

\section{问题分析}

\subsection{子问题分解}

\subsection{子问题依赖关系}

\subsection{整体技术路线}

\section{模型假设与符号说明}

\subsection{基本假设}

\subsection{符号说明}

\section{数据预处理}

\subsection{数据描述}

\subsection{缺失值处理}

\subsection{异常值剔除}

\subsection{预处理后统计描述}

\section{子问题求解}

\subsection{子问题 1: \{子问题名\} (主线)}

\subsubsection{问题描述与工程语境}
\subsubsection{模型建立}
\subsubsection{求解算法}
\subsubsection{结果与物理意义}

\subsection{子问题 2: \{子问题名\} (主线)}

\subsection{子问题 3: \{子问题名\} (主线)}

\subsection{子问题 4: \{子问题名\} (加分)}

\subsection{子问题 5: \{子问题名\} (加分)}

% ... 视题目子问数添加

\section{灵敏度与鲁棒性分析}

\subsection{工程参数扰动设计}

\subsection{联合扰动结果}

\subsection{鲁棒区间与工程含义}

\section{结果分析与方案讨论}

\subsection{分阶段方案}

\subsection{实施代价估算}

\subsection{落地风险与应对}

\section{模型评价}

\subsection{优点}

\subsection{缺点与改进}

\subsection{推广性}

% __END_SECTIONS__

\bibliographystyle{gbt7714-numerical}
\bibliography{references}

% ===== 附录 =====
\appendix

\section{主要算法伪代码}

\section{完整代码}

% TODO: 完整代码或 GitHub 链接

\section{补充图表}

\end{document}

